\section{聚类内核：源码等价流程、预算与审计}
\label{sec:clustering-implementation}

\subsection{方法选择}
\fld{method} 等于 \texttt{hybrid\_kmedoids\_tail\_5pct} 或别名
\texttt{hybrid\_kmedoids\_tail} 时进入混合路径；其他字符串使用 weighted k-medoids。
解析器不因未知字符串抛错，因此结果 audit 中的实际 \fld{clustering\_method} 是最终口径。

\subsection{候选数与代表数}
空候选直接返回空簇。实际代表数
\begin{equation}
 K'=\min(\max(1,K),N).
\end{equation}
当 $K'\ge N$ 时每个候选自成一簇，概率保持原值，距离为 0；无需进入迭代。

\subsection{连续特征列顺序}
\srcpath{feature\_summary} 从 profile 计算首值、总量、峰值、缺额和净负荷等候选摘要。
\srcpath{scaled\_features} 按确定性字段名顺序构造矩阵，避免 unordered map 遍历顺序影响距离。
开启 robust scale 时对每列应用 median/IQR；关闭时原值直接进入欧氏距离。

\subsection{风险元数据丰富}
\srcpath{raw\_risk\_source\_scores} 对 regular、reliability、resilience 使用相应风险源：负荷、
净负荷、PV/wind/renewable deficit、故障概率、停运影响等。随后
\srcpath{enrich\_candidate\_risk\_metadata}：
\begin{enumerate}
  \item 对每个源计算分位数、中位数和 IQR；
  \item 写入候选 \fld{risk\_source\_scores} 和综合 \fld{risk\_score}；
  \item 标记 global/source/coupling tail；
  \item 标记 global/source decision boundary；
  \item 汇总 \fld{anchor\_reasons} 并设置 \fld{tail\_anchor}。
\end{enumerate}

\subsection{锚点选择与截断}
\srcpath{select\_tail\_anchors} 合并风险锚点、有故障候选和 regime 极值候选并去重。若超过
$K'$，先按“含故障”布尔降序，再按 risk score 降序，最后按稳定候选顺序破同分，截断到
$K'$。audit 的 \fld{anchor\_ids} 是实际保留锚点，而非截断前全集。

\subsection{冻结语义}
仅当混合方法、\fld{include\_tail\_anchors=true} 且 \fld{freeze\_anchors=true} 时，选中锚点
成为 frozen medoid。冻结中心参与分配但不参与更新。候选本身的 \fld{tail\_anchor} 与
\fld{frozen\_medoid} 是不同字段：前者说明风险原因，后者说明优化状态。

\subsection{初始化}
锚点先进入 medoid 集；剩余位置由加权 k-means++ 型抽样填充。若总抽样权重退化为零，使用
尚未选取的确定性候选 fallback。seed 固定初始化，迭代无随机步骤。

\subsection{分配步}
每个候选计算到所有 medoid 的混合距离，选择严格更小者；完全相等时保留先出现的 medoid。
这使 medoid 排列成为破同分规则的一部分。距离同时使用连续欧氏、outage Hamming 和 regime
惩罚，签名短缺位按 0。

\subsection{更新步}
对每个非冻结非空簇，枚举簇成员 $j$，计算
\begin{equation}
 C_j=\sum_{i\in A_k}\pi_i d(i,j),
\end{equation}
选最小者作为新 medoid。空簇保持原 medoid。若所有 medoid 不变则提前停止，否则最多执行
$\max(1,\texttt{max\_iterations})$ 轮。

\subsection{簇结果}
每簇返回 cluster ID、representative ID、完整代表候选、成员概率和、成员数/ID、加权平均距离、
最大距离、冻结标志和 anchor reasons。平均距离为
\begin{equation}
 \bar d_k=\frac{\sum_{i\in A_k}\pi_i d(i,m_k)}{\Pi_k},
\end{equation}
仅在 $\Pi_k>0$ 时计算。所有簇概率最后按正总和归一。

\subsection{审计指标}
\srcpath{compute\_reduction\_metrics} 计算 transport mean/max、全局/源/耦合尾部覆盖和 regime
mass L1。\fld{compare\_baseline=true} 时，对同一候选再运行无锚点 weighted k-medoids，并保存
两组指标及 delta。它是同输入同距离的算法消融，不是与外部软件对照。

\subsection{固定 128→12 复算}
固定 seed=17、\fld{regime\_weight}=0、\fld{boundary\_fraction}=0.10 时，混合方法 12 个 medoid
全部冻结。独立 Python 从序列化候选重新计算 robust scale、成员距离、加权平均半径和全局运输
指标，最大误差为 0（容差 $10^{-10}$）。结果揭示：
\begin{center}\small
\begin{tabular}{@{}L{4.4cm}rr@{}}
\toprule
指标 & Hybrid tail-anchor & Weighted baseline\\
\midrule
transport mean & 2.148101 & 1.188380\\
transport max & 7.169293 & 3.061194\\
global tail coverage & 1.000000 & 0.000000\\
coupling tail coverage & 0.450000 & 0.150000\\
wind deficit tail coverage & 1.000000 & 0.000000\\
regime mass L1 & 1.640625 & 0.046875\\
\bottomrule
\end{tabular}
\end{center}

\subsection{选择建议}
\begin{itemize}
  \item 平均 profile 重构优先：weighted baseline 或提高 $K$；
  \item 极端后果筛查优先：hybrid，并同时报告 transport/regime 退化；
  \item 停运组合重要：提高 outage Hamming weight；
  \item regime 质量重要：提高 regime weight，但必须检查 transport 变化；
  \item 特征量纲不一致：开启 robust scale；已有物理归一化时可关闭以做审计复算。
\end{itemize}

\subsection{失败和边界}
聚类不返回“收敛到全局最优”的状态；达到迭代上限只表示停止。结果没有下游 PF/OPF 后果误差。
\begin{gapnote}
锚点数量等于 $K$ 时，所有中心被冻结，算法事实上不再优化平均距离。必须通过
\fld{frozen\_medoid\_count}、transport 和 baseline comparison 识别该边界。
\end{gapnote}
